\documentclass[10pt,a4paper]{amsart}
\usepackage[margin=19mm]{geometry}
\usepackage{amsmath,amssymb,mathtools}
\usepackage[scr=rsfs,cal=euler]{mathalfa}
\usepackage{xfrac}
\usepackage[unicode,hidelinks,bookmarks=false]{hyperref}
\newcommand{\ie}{i.e.\ }
\newcommand{\cf}{cf.\ }
\newcommand{\resp}{respectively\ }
\newcommand{\Subsection}{\subsection}
\providecommand{\nobreakdash}{\nobreak\hbox{-}}
\setlength{\emergencystretch}{2em}
\multlinegap0pt
\usepackage{bm}
\usepackage{bbm}
\usepackage{dsfont}
\usepackage{xspace}
\usepackage{cancel}
\usepackage{mathtools}
\usepackage{amsthm}
\makeatletter
\@ifundefined{theorem}{\newtheorem{theorem}{Theorem}}{}
\makeatother
\makeatletter
\newcommand{\E}{\mathbb E}
\expandafter\ifx\csname proof\endcsname\relax
\renewenvironment{proof}[1][\unskip]{%
\par
\noindent
\textbf{Proof #1.}
\noindent}
\fi
\makeatother
\title[Uncountably many conditionally inaccessible decisions exist ]{Uncountably many conditionally inaccessible decisions exist in every finite probability space}
\author{Zal\'an Gyenis, Mikl\'os R\'edei, Leszek Wro\'nski}
\date{Source: arXiv:2605.02865v1 (CC BY 4.0)}
\begin{document}
\maketitle

\noindent\emph{Source and licence.} Uncountably many conditionally inaccessible decisions exist in every finite probability space. Version arXiv:2605.02865v1 (CC BY 4.0), distributed under the Creative Commons Attribution 4.0 International licence, as recorded in the arXiv licence metadata for this exact version and shown on the abstract page https://arxiv.org/abs/2605.02865v1. Full rights notice: licenses/arxiv-2605.02865-CC-BY-4.0.txt.\par\smallskip

\noindent\textit{Editorial scope.} Counted unit: the Theorem whose source label is \texttt{thm:degrees} in the article (source0.tex lines 449--470, source label \texttt{thm:degrees}), delivered together with the article's own complete original proof of it (source0.tex lines 472--525, one \texttt{proof} environment of 54 lines). No mathematics is written, repaired or re-derived: statement and proof are the article's own text line for line. Required context retained uncounted: thm:main1 (lines 258--268); eq:def-of-r (lines 279--281); def:g (lines 322--329). Every definition, display and label of those lines is retained unchanged. Omitted from this package and not redistributed: 1--257, 269--278, 282--321, 330--448, 471--471, 526--859. Nothing inside a retained excerpt is omitted or altered: every retained body is delivered exactly as the article's lines are written, so no omission receipt of any kind accompanies this package. Citations: the retained excerpts carry no citation command at all, so no bibliography entry of the article is transcribed and no reference is redistributed. Reference adaptation: none was needed and none was made; every reference inside the delivered unit resolves inside it, so no unresolved reference or citation is produced. Printed slips kept literal: no printed word, symbol, hypothesis or number of the counted statement or of its proof was altered. Layout and class: the article's own class is \texttt{article} with packages including hyperref, amsmath, amssymb, amsopn, amsthm, enumerate, framed, xcolor, mathtools, tikz-cd, enumitem, titlesec, etoolbox, todonotes; the wrapper is the lane's pinned \texttt{amsart} editorial wrapper at 10pt on A4, which restates the theorem-like environments the retained text needs and adds the article's own macro definitions that the retained text needs (\texttt{\textbackslash E}), with extra wrapper packages (bm, bbm, dsfont, xspace, cancel, mathtools). The delivered numbering is the unit's own. Assets: no retained span contains a figure environment, a graphics-inclusion command, a PDF-page inclusion, a table or a TikZ picture; no image file is copied into this package, the package declares no asset dependency and no image file is redistributed with it. Licence: version arXiv:2605.02865v1 of this article is distributed under the Creative Commons Attribution 4.0 International licence, as recorded in the arXiv licence metadata for this exact version and shown on the abstract page https://arxiv.org/abs/2605.02865v1. A full rights notice is delivered at licenses/arxiv-2605.02865-CC-BY-4.0.txt. Characters: every retained excerpt is pure ASCII, so the delivered engine, XeLaTeX with its default Latin Modern text font, reports no missing character.
\begin{theorem}\label{thm:main1}
	For any given $p^*$ and $p$ on the measurable space $(X,\mathcal{S})$ with $X$ having $n\geq 3$ number of elements the following are equivalent:
	\begin{enumerate}
		\item[(A)] There exist real valued random variables (utility functions) $f_1$ and $f_2$ on $X$ such that the decision 
        \begin{equation}\label{eq:decision-in proof}
        \E_{p^*}[f_1]>\E_{p^*}[f_2]    
        \end{equation}  
        is conditionally $p$-inaccessible.
		\item[(B)] $p^*$ is in the Bayes $p$-Blind Spot (i.e. $q_\Pi\neq p^*$ for every proper non-trivial partition $\Pi$).
	\end{enumerate}	
\end{theorem}

	\begin{equation}\label{eq:def-of-r}
		r(i)=\frac{p^*(i)}{p(i)} \quad \text{ for } i\in X.
	\end{equation}

	\begin{equation}\label{def:g}
		g(i)=\left\{
  \begin{array}{ll}
    \log\frac{p^*(i)}{p(i)}, & \hbox{if\ } p(i)\not=0 \\
    0, & \hbox{if\ } p(i)=0
  \end{array}
\right.
	\end{equation}

\begin{theorem}\label{thm:degrees}
	Consider again the function $r$ defined by eq. (\ref{eq:def-of-r}) and assume that it is injective. Then there exists a single function $u:X\to\mathbb R$ such that the numbers
	$\{\E_q[u]:q\in\mathcal Q\}$ are pairwise distinct. Fix such a function $u$, let $g$ be the function specified by \eqref{def:g} and set $g_\eta=g+\eta u$
	for sufficiently small $\eta>0$.
	Then:
	\begin{enumerate}
		\item All values $\E_{q_\Pi}[g_\eta]$ depend only on $q_\Pi$ and are strictly smaller than $\E_{p^*}[g_\eta]$.
		\item If we list $\mathcal Q=\{q^{(1)},\dots,q^{(L)}\}$ so that
		\[
			\E_{q^{(1)}}[g_\eta]<\E_{q^{(2)}}[g_\eta]<\cdots<\E_{q^{(L)}}[g_\eta],
		\]
		and define the cumulative sums
		\[
			K_\ell=\sum_{j=1}^{\ell} m\big(q^{(j)}\big)\qquad(\ell=1,\dots,L),
		\]
		then the set of degrees realized by \emph{some} $d$ with $\E_{p^*}[d]>0$ is precisely
		\[
			\{0,\,K_1,\,K_2,\,\dots,\,K_L\}=\Big\{\sum_{j=1}^{\ell} m\big(q^{(j)}\big):\ \ell = 1, \ldots, L\Big\}.
		\]
		Moreover, each such degree is realized by a function of the form $d=g_\eta-c$ with a suitable constant $c$.
	\end{enumerate}
\end{theorem}

\begin{proof} We prove the theorem in three steps:\\
	\textbf{Step 1:} The existence of $u$.
	Think of each $q\in\mathcal Q$ as a vector in $\mathbb R^n$.
	For distinct $q\neq q'$, the set of $u\in\mathbb R^n$ with $\E_q[u]=\E_{q'}[u]$ is the hyperplane
	$\{u:\ (q-q')\cdot u=0\}$. Since there are finitely many pairs, the union of these hyperplanes is a proper subset of $\mathbb R^n$,
	so we may choose $u$ outside of it. Then $\E_q[u]$ are pairwise distinct for $q\in\mathcal Q$.

	\noindent \textbf{Step 2}: Let $g$ be as in \eqref{def:g}, and $g_\eta=g+\eta u$ for sufficiently small $\eta>0$.
	By the proof of Theorem \ref{thm:main1},
	\[
		\delta=\E_{p^*}[g]-\max_{\Pi\in\mathfrak P}\E_{q_\Pi}[g] > 0.
	\]
	Let $R$ be defined by 
    \begin{equation}
    R=\max_{q\in\mathcal Q}|\E_q[u]|+|\E_{p^*}[u]|<\infty
    \end{equation}.
	If $0<\eta<\delta/(2R)$, then
	\[
	\E_{p^*}[g_\eta]-\max_{\Pi}\E_{q_\Pi}[g_\eta]
	\ge \delta-\eta\cdot 2R>0,
	\]
	so $\E_{q_\Pi}[g_\eta]<\E_{p^*}[g_\eta]$ for all $\Pi$.

	\noindent \textbf{Step 3}:
	Because $\E_q[u]$ are pairwise distinct on $\mathcal Q$, for small $\eta$ the values
	$\E_q[g_\eta]=\E_q[g]+\eta\E_q[u]$ are also pairwise distinct on $\mathcal Q$.
	Order $\mathcal Q=\{q^{(1)},\dots,q^{(L)}\}$ increasingly by $\E_{q^{(j)}}[g_\eta]$.
	
	Fix $\ell\in\{0,1,\dots,L\}$, where $\ell=0$ means ``none''.
	Choose a constant $c$ such that
	\[
	\E_{q^{(\ell)}}[g_\eta] < c < \E_{q^{(\ell+1)}}[g_\eta],
	\]
	with the conventions $\E_{q^{(0)}}[g_\eta]=-\infty$ and $\E_{q^{(L+1)}}[g_\eta]=+\infty$.
	Define $d=g_\eta-c$. Then for a partition $\Pi$,
	\[
	\E_{q_\Pi}[d]\le 0
	\quad\Longleftrightarrow\quad
	\E_{q_\Pi}[g_\eta]\le c
	\quad\Longleftrightarrow\quad
	q_\Pi\in\{q^{(1)},\dots,q^{(\ell)}\}.
	\]
	Hence, exactly $\sum_{j\le \ell} m(q^{(j)})=K_\ell$ partitions lie in $\mathcal I(d)$, i.e.\ $\deg(d)=K_\ell$.

	Finally, since 
    \begin{equation}
    c<\E_{q^{(\ell+1)}}[g_\eta]\le \max_{\Pi}\E_{q_\Pi}[g_\eta]<\E_{p^*}[g_\eta]   
    \end{equation} we also have
	\begin{equation}
	 \E_{p^*}[d]=\E_{p^*}[g_\eta]-c>0.   
	\end{equation}
	Thus each $K_\ell$ is achievable, and no other values are possible because any $d$ induces an ordering of $\mathcal Q$
	and can only select unions of initial segments, counted with multiplicity.
\end{proof}

\end{document}
